\chapter{Анализ проблематики создания систем виртуальных мастеров для ролевых игр}
\label{chapter1}

В данной главе проводится анализ существующих подходов к созданию систем виртуальных мастеров для ролевых игр. Рассматриваются коммерческие решения, их архитектурные особенности и выявленные проблемы. Выполняется сравнительный анализ различных систем с точки зрения их способности воспроизводить функции традиционного мастера игры.

%Большие отсупы --- это хорошо. Облегчает чтение длинных <<простыней>> текста
\section{Общий анализ проблематики создания систем виртуальных мастеров}

Развитие технологий больших языковых моделей открыло новые возможности для автоматизации ролевых игр. Традиционные настольные ролевые игры требуют участия мастера игры (Game Master), который управляет миром, неигровыми персонажами и развитием сюжета. Однако создание систем, способных полноценно заменить человеческого мастера, сопряжено с рядом фундаментальных проблем.

Основная сложность заключается в том, что мастер игры выполняет множество разнородных функций: генерация повествования, управление неигровыми персонажами, контроль соблюдения правил игры, поддержание согласованности мира, балансировка сложности и обеспечение развития сюжета. Каждая из этих функций требует различных подходов к моделированию и реализации.

Современные системы, основанные на больших языковых моделях, демонстрируют высокую способность к генерации текста, однако сталкиваются с проблемами при попытке воспроизвести структурные аспекты ролевой игры. Исследования показывают, что существующие решения~\cite{Zhou2022,Jorgensen2025,Zhu2023,Santiago2023,CallisonBurch2022} в основном фокусируются на генеративной стороне процесса, но не обеспечивают необходимый уровень контроля и структурированности игрового процесса.

\section{Анализ коммерческих решений}

Наиболее известной и массовой платформой, использующей крупные языковые модели для симуляции ролевой игры в текстовом формате, является AI Dungeon~\cite{AIDungeon}. Благодаря свободной генерации текста и отсутствию жёстких ограничений по структуре сюжета, система обеспечивает высокий уровень вариативности и доступность для широкого круга пользователей. Однако её архитектура основана на чисто генеративном подходе: модель формирует ответ, опираясь исключительно на текущий контекст взаимодействия, без использования структурированного состояния игры или правил.

Это приводит к характерной особенности: система практически всегда принимает любое действие игрока, даже если оно нарушает причинно-следственные связи, изменяет факты, установленные ранее, или не согласуется с концепцией персонажа. Подтверждение любой инициативы пользователя часто приводит к разрушению логической структуры повествования и делает игровой процесс непредсказуемым в нежелательном смысле --- мир оказывается нестабильным, а сюжет лишается целостности.

Помимо AI Dungeon, появились и другие решения, стремящиеся занять нишу ИИ-мастеров~\cite{AIDungeonChi,DungeonMasterAI,HyperWriteAI}. Некоторые из них находятся в раннем доступе или работают на основе отдельных демонстрационных прототипов. Среди них можно выделить ряд платформ, предлагающих более современный интерфейс или дополнительные настройки, однако они повторяют ключевую модель взаимодействия: реактивная генеративная система, стремящаяся поддержать любые действия игрока.

Несмотря на вариации в интерфейсе и технической реализации, большинство современных альтернатив демонстрируют аналогичное поведение: ИИ не оценивает действия пользователя с точки зрения внутренней логики мира. Игровая механика и характеристики персонажей обычно представлены лишь в виде текста, а не структуры, что усиливает тенденцию к неконсистентным решениям. В ряде случаев ИИ может уходить от роли мастера и переходить в режим свободного генератора идей, что дополнительно снижает стабильность игрового процесса.

\section{Анализ проблем существующих решений}

\subsection*{Проблема чрезмерной уступчивости ИИ}

Общей чертой большинства существующих решений является неполное соответствие поведения модели ролевой функции мастера. В традиционной настольной игре мастер не только описывает мир, но и регулирует действия игроков, сохраняя внутреннюю и механическую последовательность происходящего. В текущих ИИ-системах эта функция практически отсутствует: модели ориентированы на любое поддержание диалога, что порождает чрезмерную «доброжелательность». В результате игрок может:

\begin{itemize}
	\item мгновенно получить предметы или способности, не предусмотренные историей;
	\item изменить внешний вид или факты о персонаже без последствий;
	\item выполнять действия, нарушающие физику или логические закономерности мира.
\end{itemize}

Реакция ИИ в таких случаях остаётся положительной, что лишает игру ключевых элементов структуры, напряжения и причинно-следственных связей.

\subsection*{Отсутствие слежения за состоянием мира}

Большинство доступных систем функционируют исключительно на основе текущего текста. Они не ведут явную структурированную модель мира, не хранят информацию о потребностях сюжета, состоянии мира, не отслеживают правила или ограничения, ими накладываемые. Это исключает агентность мира и целостность пространства, в котором происходит история, из-за чего не формируется общего повествования, а остаётся несвязный набор событий, происходящих последовательно.

\subsection*{Неспособность к отказу или корректирующей реакции}

В отличие от человеческого мастера, ИИ-системы не обладают механизмами обоснованного отказа. Отсутствие негативной обратной связи делает сюжет неконтролируемым, а мир --- лишённым устойчивости. Более того, эти системы не способны корректировать игрока, даже если действие явно противоречит описанным ранее фактам.

\subsection*{Потеря ролевой функции}

Некоторые доступные решения демонстрируют тенденцию к уходу от основной роли мастера. ИИ может менять рамку взаимодействия, терять «образ мастера и предлагать пользователю задачи, не связанные с игрой. Это подчёркивает отсутствие устойчивого целевого поведения и неспособность модели удерживать длительную ролевую позицию.

\section{Сравнение с традиционной ролью мастера}

Современные ИИ-мастеры в целом успешно поддерживают открытое повествование, однако им не удаётся воспроизвести структурные аспекты процесса, характерные для человеческого мастера. Среди ключевых отличий можно выделить:

\begin{itemize}
	\item неспособность к контролю повествования, что приводит к хаотичности;
	\item неумение регулировать баланс, что делает продвижение игрока чрезмерно лёгким;
	\item затруднение с удержанием долгосрочных сюжетов, поскольку контекст моделируется только текстом.
\end{itemize}

Эти особенности ограничивают применимость существующих решений в качестве полноценных инструментов для проведения настольных ролевых кампаний и подчеркивают их преимущественно развлекательный характер.






\section{Сравнительный анализ программных средств}

Для систематизации анализа существующих решений были рассмотрены основные коммерческие платформы и выделены их ключевые характеристики.

Анализ показал, что все рассмотренные коммерческие системы демонстрируют схожие ограничения: отсутствуют механизмы контроля действий игрока, не используется структурированное представление состояния мира, отсутствует интеграция механических правил игры, а также не предусмотрена система управления неигровыми персонажами с учётом их характеристик и эмоционального состояния.

Исследовательские работы~\cite{Zhou2022,Zhu2023,CallisonBurch2022} предлагают более продвинутые подходы, включающие использование теории разума (theory of mind) для персонажей, системы намерений и более структурированное управление диалогом. Однако эти решения остаются на стадии прототипов и не интегрируют все необходимые компоненты в единую систему.




\section{Выводы}

На основе проведённого анализа существующих решений для создания систем виртуальных мастеров ролевых игр можно сделать следующие выводы:

\begin{enumerate}
	\item Выполнен анализ коммерческих платформ (AI Dungeon, AI Dungeon-Chi, Dungeon Master AI, HyperWrite AI) и исследовательских работ в области применения больших языковых моделей для ролевых игр. Все рассмотренные коммерческие решения основаны на чисто генеративном подходе без структурированного управления состоянием игры.
	
	\item Выявлены ключевые проблемы существующих систем: чрезмерная уступчивость ИИ, отсутствие слежения за состоянием мира, неспособность к отказу или корректирующей реакции, потеря ролевой функции мастера. Эти проблемы ограничивают применимость систем в качестве полноценных инструментов для проведения ролевых кампаний.
	
	\item Сравнительный анализ показал, что ни одна из существующих систем не обеспечивает необходимый уровень контроля действий игрока, структурированного представления состояния мира, интеграции механических правил и управления неигровыми персонажами. Требуется разработка новой архитектуры, объединяющей эти компоненты.
	
	\item Исследовательские работы демонстрируют перспективные направления, такие как использование теории разума для персонажей и систем намерений, однако они не интегрируют все необходимые компоненты в единую систему. На основе анализа принято решение разработать систему, включающую: структурированное представление состояния игры, механизмы контроля действий игрока, систему управления эмоциями и намерениями персонажей, интеграцию механических правил игры.
\end{enumerate}



\section{Постановка задачи курсовой работы}

На основе проведённого анализа существующих решений и выявленных проблем сформулирована цель работы: разработка системы ролевой игры с виртуальным мастером, которая обеспечивает структурированное управление игровым процессом, контроль действий игрока и реалистичное поведение неигровых персонажей на основе их эмоционального состояния и намерений.

Для достижения поставленной цели необходимо решить следующие задачи:

\begin{enumerate}
	\item Разработать архитектуру системы ролевой игры, включающую компоненты для управления игровым циклом, генерации повествования, управления персонажами и контроля согласованности действий.
	
	\item Реализовать систему управления эмоциональным состоянием и намерениями персонажей на основе моральных схем, обеспечивающую динамическое изменение поведения персонажей в зависимости от их взаимодействий.
	
	\item Разработать механизм директора (Director), который анализирует текущее состояние игры и выдаёт директивы неигровым персонажам для продвижения сюжета, обеспечивая баланс между свободой действий игрока и следованием задуманному сюжету.
	
	\item Реализовать систему сторителлинга (Storyteller), генерирующую связное повествование на основе действий персонажей и текущего состояния игры.
	
	\item Интегрировать механические элементы игры (броски кубиков, проверки характеристик) с системой генерации повествования для обеспечения согласованности между механическими правилами и описанием происходящего.
	
	\item Реализовать механизмы контроля действий игрока и неигровых персонажей на соответствие предыстории персонажей, логике мира и правилам игры.
	
	\item Провести экспериментальную проверку разработанной системы и сравнение с существующими аналогами.
\end{enumerate}

Разрабатываемая система должна обеспечивать более структурированный и контролируемый игровой процесс по сравнению с существующими генеративными решениями, сохраняя при этом гибкость и вариативность повествования, характерные для традиционных ролевых игр с живым мастером. 

%%% Local Variables:
%%% TeX-engine: xetex
%%% eval: (setq-local TeX-master (concat "../" (seq-find (-cut string-match ".*-3-pz\.tex$" <>) (directory-files ".."))))
%%% End:
